\par\bigskip\noindent\textbf{K101 supporting proof: An independent actual selected-corner atlas}\par\nobreak
Working exact-map audit, owned by blind\_control. The source, Cartier
and root inputs are EI.1-\/-9 and EIR.1-\/-10. This note constructs the
selected \(T_0\) toric patch and its source-graph refinements, with
the full period and line maps. It assigns no characteristic number
or source return. The original finite cap stays on the domain in
EIR.9 and retains the EIR.28 annular map.

\par

\medskip\noindent\textbf{SIT.1. Actual periods at the four selected branches}

\par\nobreak

On the ordinary chart put
\[
 e=\chi\sigma,\qquad
 H=1-\chi^4(1-\chi^2\sigma^2).
 \tag{SIT.1}
\]
At \((\chi,\sigma)=(\chi_0,0)\), \(\chi_0^4=1\), the
\(\chi\)-derivative of \(H\) is \(-4\chi_0^3\). Thus the actual
inverse function gives \(\chi=\chi(\sigma,H)\), with no branch
exchange on its chosen small disc. The parameter \(e=\chi\sigma\)
retains the nonzero unit \(\chi\). The source Tate coordinate is
\[
 q=Q=\kappa(\sigma,H)H ,
 \tag{SIT.2}
\]
where \(\kappa\) is the full analytic nonzero height unit. The point
line coordinate is \(V=f_{P^a}(Z)v_M\), with the entire theta
function, and the actual source translation is \(Z\mapsto gZ\),
\(g=P^{-b}\). Its full line coefficient is
\[
 \beta=\sigma^{-6}\bar B_{a,b}(p)
       \frac{f_{P^a}(gZ)}{f_{P^a}(Z)} .
 \tag{SIT.3}
\]
The horizontal divisors cancel by EIR.47. Both numerator and
denominator retain their actual Cartier arrows. The result is a
nonzero holomorphic function on each complete elliptic source
fibre, and is therefore constant on that fibre. Evaluation at a
fixed regular Tate point gives its analytic extension in the
base; at \(\sigma=0\), \(P=-1\), \(a\) is even and the complete
coefficient is \(64\). This proves the exact base unit, with its
source/target direction fixed by (SIT.3).

Consequently the periods are precisely
\[
 h:(Z,V)\longmapsto(qZ,\alpha V),\quad \alpha=P^a,
\qquad
 B:(Z,V)\longmapsto(gZ,\sigma^4q\,C V),\quad
 C=\lambda\beta/\kappa .
 \tag{SIT.4}
\]
No part of \(\kappa,\alpha,g,C\) is replaced by its central value.
The valuation matrices in the physical order \((Z,V)\), for the
base divisors \((\sigma,H)\), are
\[
 E=\begin{pmatrix}0&0\\0&4\end{pmatrix},
 \qquad I=\begin{pmatrix}1&0\\0&1\end{pmatrix}.
 \tag{SIT.5}
\]
The remaining \(\mu_4\) fixes \(e,H\) and the native source point,
and permutes the four \(\chi_0\) discs freely. In the displayed
source/point frame it fixes \(Z,V,q\) and sends \(\sigma\) to
\(-i\sigma\). Its full finite arrow is retained below.

\par

\medskip\noindent\textbf{SIT.2. Six regular charts and the exact negative boundary fan}

\par\nobreak

On the source chain use
\[
 x_i y_i=q,\qquad Z=q^i x_i.
\]
For the line chain use
\[
 U_j W_j=\sigma^4q,\qquad V=(\sigma^4q)^j U_j.
 \tag{SIT.6}
\]
Resolve their combined equation with the ordered factors
\[
 (f_1,\ldots,f_6)=(\sigma,\sigma,\sigma,\sigma,y_i,x_i).
 \tag{SIT.7}
\]
The \(r\)-th actual chart is
\[
 f_r=s_rt_r,\quad
 U_j=s_r\prod_{\ell<r}f_\ell,\quad
 W_j=t_r\prod_{\ell>r}f_\ell .
 \tag{SIT.8}
\]
It is smooth: \(f_r\) is one of the independent source-base
coordinates \(\sigma,x_i,y_i\), and replacing it by \(s_rt_r\)
leaves four independent coordinates. On an adjacent overlap,
\[
 s_{r+1}=t_r^{-1},\qquad t_{r+1}=f_{r+1}t_r .
 \tag{SIT.9}
\]
These maps have the written inverse and preserve both (SIT.6)
and every remaining factor.

In the original lattice
\(N=(\operatorname{ord}Z,\operatorname{ord}V,
\operatorname{ord}\sigma,\operatorname{ord}H)\), put
\[
 E_j=(0,j,1,0),\qquad I_n=(n_1,n_2,0,1).
\]
The six maximal cones are exactly
\[
 \begin{array}{ll}
 (E_{r-1},E_r,I_0,I_{e_1}),&1\le r\le4,\\
 (E_4,I_0,I_{e_1},I_{e_2}),&\\
 (E_4,I_{e_1},I_{e_2},I_{e_1+e_2}).&
 \end{array}
 \tag{SIT.10}
\]
Each determinant has absolute value one. The last two triangles
on the \(\sigma\)-unit boundary are
\((0,e_1,e_2)\) and \((e_1,e_2,e_1+e_2)\). Thus they give the
current negative fan \(T_0\), with its diagonal \(e_1-e_2\).

The entire period collection has exact common faces and support.
At fixed positive base heights \((R,S)\), the first four cones
cut out the rectangles
\[
 [0,S]\times[(r-1)R,rR]\quad(1\le r\le4).
\]
The last two cones cut out the square
\([0,S]\times[4R,4R+S]\), split by its negative diagonal.
The actual periods are \((S,0)\) and \((0,4R+S)\). Their
translates tile the plane by these whole rectangles and triangles.
Both the right/left edges and top/bottom edges match by these
literal period vectors. The resulting four-dimensional cones
therefore meet along entire common faces; homogeneity gives
the same assertion at every height ratio. The boundary at
\(S=0\) retains every primitive chain ray \(E_j\).

\par

\medskip\noindent\textbf{SIT.3. Full units, both chain overlaps and the finite action}

\par\nobreak

For a simultaneous factor change \(f_\ell'=a_\ell f_\ell\)
and \(U_j'=cU_j\), with \(\prod a_\ell=1\), the actual lift is
\[
 s_r'=\frac{c}{\prod_{\ell<r}a_\ell}s_r,\qquad
 t_r'=\frac{a_r\prod_{\ell<r}a_\ell}{c}t_r .
 \tag{SIT.11}
\]
This is an equality of analytic maps, retaining all parameter
derivatives of the nonzero units. For \(h\) the labels are
\((i,j)\mapsto(i+1,j)\), all \(a_\ell=1\), and \(c=\alpha\).
For \(B\) the labels are \((i,j)\mapsto(i,j+1)\),
\[
 (a_1,\ldots,a_6)=(1,1,1,1,g^{-1},g),\qquad c=C.
 \tag{SIT.12}
\]
Substitution in (SIT.8) gives exactly (SIT.4).
These lifts commute, because their base units are fixed by
both source translations and their dense-torus maps commute.
Their equality holds on the entire charts by the explicit
formulas, including the normal derivatives along \(q=0\).

The source-chain overlap is
\[
 x_{i+1}=y_i^{-1},\qquad y_{i+1}=x_i y_i^2.
 \tag{SIT.13}
\]
On \(y_i\ne0\), the old last-factor chart
\(x_i=s_6t_6\) maps to the new penultimate-factor chart by
\[
 s_5'=y_i s_6,\qquad t_5'=y_i t_6,\qquad x_{i+1}=y_i^{-1}.
 \tag{SIT.14}
\]
Indeed \(U_j=\sigma^4y_i s_6=\sigma^4s_5'\) and
\(W_j=t_6=t_5'x_{i+1}\). In the first four charts \(s_r,t_r\)
are unchanged and (SIT.13) is substituted in the source
coordinates. Unit-factor deletion identifies the other
overlaps with these same maps.

On \(W_j\ne0\), the line-chain overlap is
\[
 U_{j+1}=W_j^{-1},\qquad W_{j+1}=\sigma^4q\,W_j.
 \tag{SIT.15}
\]
In the last chart \(W_j=t_6\) this maps to the next first chart
by \(s_1'=t_6^{-1}\), \(t_1'=\sigma t_6\), with the same
source coordinates. The inverse is obtained from these two
displayed equalities. Both physical coordinates \(Z,V\) are
therefore identical across all source and line overlaps.

On the \(q\)-unit overlap with the ordinary chain of EIR.10,
delete the nonzero source factors in (SIT.7). For the remaining
\(r\)-th \(\sigma\) chart, \(1\le r\le4\), the exact arrow is
\[
 k=4j+r-1,\qquad X_k=q^j s_r,\qquad Y_k=q^{-j}t_r.
 \tag{SIT.16}
\]
It follows by substituting
\(V=(\sigma^4q)^j\sigma^{r-1}s_r\).
The \(B\)-multiplier becomes \(qC=\lambda H\beta\), as
required in the same point-theta frame. On returning to the
genuine point frame, the full theta/Cartier ratio in (SIT.3)
gives exactly \(\lambda H\mathcal B_{a,b}\) in EIR.53.
More explicitly, let \(r_a\) be the genuine point Cartier
function in EIR.48 and put
\(\psi=f_{P^a}/r_a\). Its complete numerator and denominator
have the same Cartier divisor, including both total transforms
on the source graph, so \(\psi\) and its reciprocal are
holomorphic units on the specified overlap. The point-theta
coordinate is \(\psi\) times the genuine point coefficient.
The latter chain therefore has the actual map
\[
 X_k=\psi^{-1}q^js_r,\qquad
 Y_k=\psi q^{-j}t_r,\qquad k=4j+r-1.
 \tag{SIT.16a}
\]
The ratio of the target and source \(\psi\)\textquotesingle s changes
\(\beta\) to \(\mathcal B_{a,b}\) exactly. In particular, both
the factor \(q^j\) and the entire point-line unit are retained.

Finally the remaining finite arrow has
\[
 a_1=\cdots=a_4=-i,\quad a_5=a_6=1,\quad c=1.
 \tag{SIT.17}
\]
Formula (SIT.11) gives its complete lift. On the first four
charts it is \(s_r\mapsto i^{r-1}s_r\),
\(t_r\mapsto(-i)^r t_r\), and it fixes \(s_r,t_r\) on
the last two. It also permutes the four selected source
branches. The fourth iterate is the identity. This is the
actual finite source/line map on the selected atlas.

\par

\medskip\noindent\textbf{SIT.4. The source O/T graph and its relative nine-factor resolution}

\par\nobreak

In the \(O\) source chart let
\(\eta=-\bar X/\bar Y\), \(U=\eta^2\bar X\). Its equation is
\[
 1=U+a_2\eta^2+d\eta^4/U+\sigma^2\eta^6/U^2 .
 \tag{SIT.18}
\]
The source graph blowup chart \(\sigma=\eta v\) is smooth,
with \(U\) an analytic unit. The actual \(-G\) chord gives
\[
 \begin{aligned}
 k&=-\eta^{-1}+\eta^3v/U,\\
 \bar X'&=\eta^2 A_1,\qquad
 A_1=(d-2v)/U+2\eta^4v^2/U^2,\\
 \bar Y'&=\eta c,\qquad
 c=v+A_1-\eta^4vA_1/U.
 \end{aligned}
 \tag{SIT.19}
\]
The target \(T\)-graph is \(\sigma=\bar Y'v'\), so the full
source arrow is
\[
 (\eta,v)\longmapsto(\eta c,v/c).
 \tag{SIT.20}
\]
Here \(c|_{\eta=0}=1-v\). It is a unit on the chosen corner
neighbourhood. The inverse chord gives the inverse analytic
map. For \(-3G\), compose the same three chord maps; its
result is again \((\eta,v)\mapsto(\eta c_3,v/c_3)\), with
the complete composed unit retained. The central source
coordinates satisfy \(p=\eta/\sigma=1/v\) and
\(z=\bar Y'/\sigma=p-1\); successive translations give
\(p\mapsto p-b\), so \(c_b|_{\eta=0}=1-bv\), for
both admitted \(b=1,3\).

The pullback of (SIT.6) to this genuine source graph is
\[
 U_jW_j=\eta^4v^4q.
 \tag{SIT.21}
\]
Its relative ordered resolution has factors
\[
 (\eta,\eta,\eta,\eta,v,v,v,v,q).
 \tag{SIT.22}
\]
Formula (SIT.8), now with nine factors, gives nine smooth
four-dimensional charts. Every factor is one of the three
independent local coordinates \(\eta,v,q\); its monomial
matrix has determinant one in absolute value. Formula
(SIT.11) lifts the actual source/line map with
\[
 (a_1,\ldots,a_9)=
 (c_b,c_b,c_b,c_b,c_b^{-1},c_b^{-1},c_b^{-1},c_b^{-1},1).
 \tag{SIT.23}
\]
The line multiplier \(c\) in (SIT.11) is the full Hom
coefficient from (SIT.4), pulled to this source graph.
It is distinct from the source unit \(c_b\).
Thus the source translation preserves the two retained
factor divisors and has a literal lift through the resolution.

On \(\eta\ne0\), delete its four unit factors and use
\(\sigma=\eta v\). On \(v\ne0\), delete its four unit factors
and use the same equality. Formula (SIT.11) gives both
resulting unit rescalings, preserving the entire \(U_j,W_j\)
and \(q\). These are the actual overlaps with the four-\(\sigma\)
factor atlas. The remaining finite map has
\(\eta\mapsto-i\eta\), \(v\mapsto v\); its first four
factor multipliers are \(-i\), and all the rest are one.
The identical construction applies in the \(T\) graph
\(\sigma=\bar Yv\), with the full \(O/T\) chord square (SIT.20).

This completes the written local \(T_0\) selected and O/T corner
atlas. Its pieces retain the ordinary line-chain overlap,
the source graph, the finite action and the old-cap comparison.
The compact global quotient, all its boundary components and
any associated characteristic evaluation still require their
joint assembly; no count is extracted from this atlas alone.

\par

\medskip\noindent\textbf{SIT.5. Remaining finite normal maps at the ordinary section corner}

\par\nobreak

The source graph must be carried before taking its finite quotient.
On the ordinary chain its exact corner equation and fourth-root
weights are
\[
 X_jY_j=\eta v,\qquad
 (\chi,\eta,v,X_j,Y_j)\longmapsto
 (i\chi,-i\eta,v,i^jX_j,i^{-j-1}Y_j).
 \tag{SIT.24}
\]
Use the actual two-factor resolution in the order \((\eta,v)\).
Its charts and complete complex-linear normal weights modulo four
are
\[
 \begin{array}{c|c|c}
 \text{chart}&\text{coordinate equations}&\mu_4\text{ weights}\\\hline
 1&\eta=st,\ X_j=s,\ Y_j=tv
       &(\chi,v,s,t):(1,0,j,-j-1)\\
 2&v=st,\ X_j=\eta s,\ Y_j=t
       &(\chi,\eta,s,t):(1,-1,j+1,-j-1).
 \end{array}
 \tag{SIT.25}
\]
Both rows follow by direct substitution in (SIT.24); their
determinant character is one. In both rows the physical circle
has the actual weights \((0,0,+1,-1)\), in the displayed
coordinate order. All four index types remain:
\[
 \begin{array}{c|c|c}
 j\bmod4&\text{chart 1}&\text{chart 2}\\\hline
 0&(1,0,0,3)&(1,3,1,3)\\
 1&(1,0,1,2)&(1,3,2,2)\\
 2&(1,0,2,1)&(1,3,3,1)\\
 3&(1,0,3,0)&(1,3,0,0).
 \end{array}
 \tag{SIT.26}
\]
In particular, cover smoothness does not identify the finite
coarse quotient with that cover: the full generator has isolated
fixed points in the second chart when \(j=0,1,2\), and its
linear action is the displayed nonidentity action. The source
reset carries \(j\) to \(j+4\), exchanges the O/T corner, and
has the actual invariant unit \(c_b\) of (SIT.20). Hence it
preserves these complete finite normal representations and
their line maps. They are retained for the remaining finite
quotient resolution and boundary calculation.

